\begin{prob}
    Consider the algorithm below.

    \begin{minted}[autogobble]{python}
        def bogosearch(numbers, target):
            """search by randomly guessing. `numbers` is an array of n numbers"""
            n = len(numbers)

            while True:
                # randomly choose a number between 0 and n-1 in constant time
                guess = np.random.randint(n)
                if numbers[guess] == target:
                    return guess
    \end{minted}

    We will set up the analysis of the expected time complexity of this algorithm.

    \begin{subprobset}
        \begin{subprob}
            What are the cases? How many are there?

            \vspace{3em}

            \begin{soln}
                Case $\alpha$ occurs when the target is found on iteration $\alpha$.
                In principle, the algorithm can run forever, although this is very
                unlikely (it happens with probability zero). As such, there are
                infinitely-many cases.
            \end{soln}
        \end{subprob}

        \begin{subprob}
            What is the probability of case $\alpha$?

            \vspace{3em}

            \begin{soln}
                $P(\alpha)$ is the probability of guessing wrong $\alpha - 1$
                times and guessing right on the $\alpha$th time:
                \[
                    (1 - 1/n)^{\alpha - 1} \cdot (1/n)
                \]
            \end{soln}
        \end{subprob}

        \begin{subprob}
            What is the running time in case $\alpha$?

            \vspace{3em}

            \begin{soln}
                We perform $\alpha$ iterations in case $\alpha$, each taking
                constant time, $c$. The total work in case $\alpha$ is therefore
                $c \alpha$.
            \end{soln}
        \end{subprob}
    \end{subprobset}

\end{prob}
